% ROZDZIAŁ 1 — plik włączany przez main.tex
\chapter{Rodziny zbiorów i przestrzenie mierzalne}
% Konwencja obowiązująca w całym skrypcie: funkcję charakterystyczną zbioru A oznaczamy przez $\chi_A$.
\label{ch:przestrzenie-mierzalne}

\section{Dlaczego potrzebujemy \texorpdfstring{$\sigma$}{sigma}-algebr?}

W teorii miary chcemy przypisywać rozmiar zbiorom, a następnie całkować funkcje.
Aksjomaty teorii mają zapewniać zgodność tego rozmiaru z rozkładaniem zbioru
na przeliczalnie wiele rozłącznych części. Nie wystarczy zatem umieć mierzyć
pojedynczych zbiorów: potrzebujemy rodziny zbiorów zamkniętej na operacje,
które pojawiają się w granicach ciągów i szeregach. Właśnie tę rolę pełnią
$\sigma$-algebry. Istotne jest, że nie zakładamy mierzalności wszystkich
podzbiorów rozważanej przestrzeni.

W całym rozdziale $X$ i $Y$ oznaczają zbiory, $\mathcal P(X)$ --- zbiór
wszystkich podzbiorów $X$, natomiast $A^c=X\setminus A$ oznacza dopełnienie
zbioru $A\subseteq X$ względem ustalonej przestrzeni $X$. Symbol $\mathbb N$
oznacza zbiór liczb naturalnych dodatnich. Rodziny zbiorów oznaczamy
$\mathcal A,\mathcal F,\mathcal G,\mathcal M$.

\begin{definicja}[Pierścień, algebra, $\sigma$-algebra]
\label{def:algebry}
Niech $X$ będzie zbiorem, a $\mathcal R\subseteq\mathcal P(X)$.
\begin{enumerate}
\item $\mathcal R$ nazywamy \emph{pierścieniem zbiorów}, jeżeli
$\varnothing\in\mathcal R$ i dla dowolnych $A,B\in\mathcal R$ zachodzą
$A\cup B\in\mathcal R$ oraz $A\setminus B\in\mathcal R$.
\item $\mathcal R$ nazywamy \emph{algebrą zbiorów} na $X$, jeżeli jest
pierścieniem i $X\in\mathcal R$.
\item $\mathcal R$ nazywamy \emph{$\sigma$-algebrą} na $X$, jeżeli
$X\in\mathcal R$, $A^c\in\mathcal R$ dla każdego $A\in\mathcal R$ oraz
$\bigcup_{n=1}^{\infty}A_n\in\mathcal R$ dla każdego ciągu
$(A_n)_{n\geq1}$ elementów $\mathcal R$.
\end{enumerate}
Parę $(X,\mathcal F)$, gdzie $\mathcal F$ jest $\sigma$-algebrą na $X$,
nazywamy \emph{przestrzenią mierzalną}; jej elementy to \emph{zbiory
mierzalne}.
\end{definicja}

\begin{uwaga}
Algebra jest zamknięta na skończone sumy, przekroje i różnice zbiorów.
$\sigma$-algebra jest ponadto zamknięta na przeliczalne przekroje, gdyż
z praw de Morgana
$\bigcap_{n\geq1}A_n=(\bigcup_{n\geq1}A_n^c)^c$.
Każda $\sigma$-algebra jest algebrą, ale nie odwrotnie.
\end{uwaga}

\begin{przyklad}[Najmniejsza i największa $\sigma$-algebra]
Dla dowolnego $X$ rodziny $\{\varnothing,X\}$ oraz $\mathcal P(X)$ są
$\sigma$-algebrami. Pierwszą nazywamy \emph{trywialną}, drugą ---
\emph{dyskretną}. Gdy $X=\varnothing$, są one równe.
\end{przyklad}

\begin{przyklad}[Zbiory skończone i kofinitywne]
\label{ex:finite-cofinite}
Niech $X$ będzie nieskończonym zbiorem. Rodzina
\[
\mathcal A=\{A\subseteq X:A\ \text{jest skończony lub }X\setminus A
\ \text{jest skończony}\}
\]
jest algebrą. Jeśli $X=\mathbb N$, nie jest $\sigma$-algebrą:
każdy singleton $\{2n\}$ należy do $\mathcal A$, lecz
$\bigcup_{n\geq1}\{2n\}$ ani nie jest skończony, ani nie ma skończonego
dopełnienia.
\end{przyklad}

\begin{przyklad}[Algebra przeliczalno-koprzeliczalna]
\label{ex:countable-cocountable}
Niech $X$ będzie dowolnym zbiorem. Rodzina
\[
\mathcal C_X=\{A\subseteq X:A\text{ jest przeliczalny lub }X\setminus A
\text{ jest przeliczalny}\}
\]
jest $\sigma$-algebrą. Istotnie, dopełnienia zachowują podaną własność.
Jeżeli w przeliczalnej rodzinie $(A_n)$ występuje zbiór o przeliczalnym
dopełnieniu, to dopełnienie $\bigcup_nA_n$ jest przeliczalne; jeśli zaś
wszystkie $A_n$ są przeliczalne, ich suma również jest przeliczalna.
Jeżeli $X$ jest nieprzeliczalny, to $\mathcal C_X\subsetneq\mathcal P(X)$.
Rzeczywiście, możemy wybrać dwa rozłączne nieprzeliczalne podzbiory
$A,B$ o sumie $X$ (wybór takiego podziału jest możliwy w zwykłych
założeniach teorii mnogości z aksjomatem wyboru). Wtedy ani $A$,
ani $X\setminus A=B$ nie jest przeliczalny, więc
$A\notin\mathcal C_X$.
\end{przyklad}

\begin{przyklad}[Zbiory określone przez podział]
Niech $(E_i)_{i\in I}$ będzie rozbiciem zbioru $X$ na niepuste, parami
rozłączne zbiory. Rodzina wszystkich sum $\bigcup_{i\in J}E_i$ dla
$J\subseteq I$ jest $\sigma$-algebrą. W szczególności dla
$X=\{a,b,c\}$ podział $\{\{a,b\},\{c\}\}$ wyznacza $\sigma$-algebrę
$\{\varnothing,\{a,b\},\{c\},X\}$.
\end{przyklad}

\section{Generowanie \texorpdfstring{$\sigma$}{sigma}-algebr}

\begin{twierdzenie}[Przekroje $\sigma$-algebr]
\label{thm:intersection-sigma}
Niech $X$ będzie zbiorem i niech $(\mathcal F_i)_{i\in I}$ będzie
niepustą rodziną $\sigma$-algebr na $X$. Wówczas
$\bigcap_{i\in I}\mathcal F_i$ jest $\sigma$-algebrą na $X$.
\end{twierdzenie}
\begin{proof}
Wszystkie $\mathcal F_i$ zawierają $X$. Jeśli $A$ należy do każdego
$\mathcal F_i$, to jego dopełnienie również należy do każdego z nich.
Analogicznie, jeśli wszystkie $A_n$ należą do przekroju tych rodzin,
to $\bigcup_nA_n$ należy do każdej $\mathcal F_i$, a więc do ich przekroju.
\end{proof}

\begin{definicja}[$\sigma$-algebra generowana]
\label{def:sigma-generated}
Dla zbioru $X$ i rodziny $\mathcal E\subseteq\mathcal P(X)$ definiujemy
\[
\sigma_X(\mathcal E)=\bigcap\{\mathcal F:\mathcal F\text{ jest
$\sigma$-algebrą na }X\text{ i }\mathcal E\subseteq\mathcal F\}.
\]
Jest to \emph{$\sigma$-algebra generowana przez $\mathcal E$}.
Jeśli przestrzeń $X$ jest ustalona, piszemy po prostu $\sigma(\mathcal E)$.
\end{definicja}

Rodzina pod przekrojem jest niepusta, ponieważ zawiera $\mathcal P(X)$.
Z twierdzenia \ref{thm:intersection-sigma} wynika, że $\sigma(\mathcal E)$
jest $\sigma$-algebrą. Ponadto jest najmniejszą $\sigma$-algebrą
zawierającą $\mathcal E$, w sensie inkluzji.

\begin{przyklad}
Dla $E\subseteq X$ mamy
$\sigma(\{E\})=\{\varnothing,E,E^c,X\}$, z ewentualnymi powtórzeniami,
gdy $E=\varnothing$ lub $E=X$.
\end{przyklad}

\begin{uwaga}[Pułapka dotycząca sumy rodzin]
Nawet rosnąca suma $\sigma$-algebr nie musi być $\sigma$-algebrą.
Na $X=\mathbb N$ niech $\mathcal F_n$ oznacza $\sigma$-algebrę
generowaną przez singletony $\{1\},\ldots,\{n\}$.
Wówczas $\mathcal F_n\subseteq\mathcal F_{n+1}$, ale
$\bigcup_n\mathcal F_n$ jest tylko algebrą zbiorów skończonych i
kofinitywnych, znaną z przykładu \ref{ex:finite-cofinite}.
\end{uwaga}

\begin{definicja}[$\sigma$-algebra śladowa]\label{idx:01:definicja:04}

Niech $(X,\mathcal F)$ będzie przestrzenią mierzalną i niech
$E\subseteq X$ będzie dowolnym podzbiorem. Definiujemy
$\mathcal F|_E=\{A\cap E:A\in\mathcal F\}$.
Rodzina ta jest $\sigma$-algebrą na $E$, nazywaną \emph{śladem}
$\mathcal F$ na $E$.
\end{definicja}
\begin{proof}[Uzasadnienie]
Do rodziny należy $E=X\cap E$. Dopełnienie w $E$ zbioru $A\cap E$
wynosi $A^c\cap E$, a przeliczalna suma zbiorów $A_n\cap E$ jest równa
$(\bigcup_n A_n)\cap E$.
\end{proof}

\section{Układy \texorpdfstring{$\pi$}{pi} i układy Dynkina}

\begin{definicja}[Układ $\pi$ i układ Dynkina]
\label{def:pi-lambda}
Niech $X$ będzie zbiorem. Rodzinę $\mathcal P\subseteq\mathcal P(X)$
nazywamy \emph{układem $\pi$}, jeżeli jest niepusta i
$A\cap B\in\mathcal P$ dla wszystkich $A,B\in\mathcal P$.
Rodzinę $\mathcal D\subseteq\mathcal P(X)$ nazywamy \emph{układem
Dynkina} (lub \emph{układem $\lambda$}), jeżeli:
\begin{enumerate}
\item $X\in\mathcal D$;
\item jeśli $A,B\in\mathcal D$ i $A\subseteq B$, to $B\setminus A\in\mathcal D$;
\item dla dowolnych parami rozłącznych $A_1,A_2,\ldots\in\mathcal D$
zachodzi $\bigcup_nA_n\in\mathcal D$.
\end{enumerate}
\end{definicja}

Każda $\sigma$-algebra jest układem Dynkina i układem $\pi$.
Samo bycie układem Dynkina nie gwarantuje zamknięcia na przekroje.
Przykładowo, jeśli $X=\{1,2,3,4\}$, to rodzina
\[
\mathcal D=\{\varnothing,X,\{1,2\},\{3,4\},\{1,3\},\{2,4\}\}
\]
jest układem Dynkina (jedynymi właściwymi, niepustymi zbiorami są tu
zbiory dwuelementowe), ale nie jest $\sigma$-algebrą, gdyż
$\{1,2\}\cap\{1,3\}=\{1\}\notin\mathcal D$.

\begin{lemat}
\label{lem:dynkin-intersection}
Układ Dynkina $\mathcal D$ na zbiorze $X$ jest $\sigma$-algebrą wtedy
i tylko wtedy, gdy jest zamknięty na skończone przekroje.
\end{lemat}
\begin{proof}
Konieczność jest oczywista. Załóżmy, że $\mathcal D$ jest układem
Dynkina zamkniętym na przekroje. Z warunku drugiego otrzymujemy
$A^c=X\setminus A\in\mathcal D$. Dla $A,B\in\mathcal D$ zachodzi
$B\setminus A=B\setminus(A\cap B)\in\mathcal D$, więc
$A\cup B=A\,\dot\cup\,(B\setminus A)\in\mathcal D$; symbol
$\dot\cup$ oznacza sumę rozłączną. Powtarzając ten krok, uzyskujemy
zamknięcie na sumy skończone. Dla dowolnego ciągu $(A_n)$ elementów
$\mathcal D$ zastępujemy $A_n$ przez
$B_n=A_n\setminus\bigcup_{k<n}A_k$. Zbiory $B_n$ są mierzalne względem
$\mathcal D$, rozłączne, a ich suma jest równa $\bigcup_nA_n$.
Trzeci aksjomat układu Dynkina kończy dowód.
\end{proof}

\begin{twierdzenie}[Twierdzenie $\pi$--$\lambda$ Dynkina]
\label{thm:pi-lambda}
Niech $X$ będzie zbiorem, $\mathcal P$ układem $\pi$ na $X$,
a $\mathcal D$ układem Dynkina na $X$. Jeśli
$\mathcal P\subseteq\mathcal D$, to
$\sigma(\mathcal P)\subseteq\mathcal D$.
\end{twierdzenie}
\begin{proof}
Przekrój dowolnej rodziny układów Dynkina na $X$ jest układem Dynkina.
Niech $\mathcal L$ będzie najmniejszym układem Dynkina zawierającym
$\mathcal P$. Wystarczy pokazać, że $\mathcal L$ jest $\sigma$-algebrą.

Dla ustalonego $A\in\mathcal P$ definiujemy
$\mathcal L_A=\{B\in\mathcal L:A\cap B\in\mathcal L\}$.
Rodzina $\mathcal L_A$ jest układem Dynkina: należy do niej $X$, gdyż
$A\cap X=A\in\mathcal L$; własność różnicy wynika z równości
$A\cap(B\setminus C)=(A\cap B)\setminus(A\cap C)$ dla $C\subseteq B$;
własność rozłącznych sum wynika z rozdzielności przekroju względem sumy.
Ponieważ $\mathcal P$ jest układem $\pi$, mamy
$\mathcal P\subseteq\mathcal L_A$. Z minimalności $\mathcal L$ wynika
$\mathcal L_A=\mathcal L$. Oznacza to, że
$A\cap B\in\mathcal L$ dla wszystkich $A\in\mathcal P$, $B\in\mathcal L$.

Teraz dla ustalonego $B\in\mathcal L$ rozważamy
$\mathcal K_B=\{A\in\mathcal L:A\cap B\in\mathcal L\}$.
Ponownie jest to układ Dynkina. Wykazana własność daje
$\mathcal P\subseteq\mathcal K_B$, skąd
$\mathcal K_B=\mathcal L$. Zatem $\mathcal L$ jest zamknięty
na skończone przekroje. Na mocy lematu \ref{lem:dynkin-intersection}
jest $\sigma$-algebrą, a więc $\sigma(\mathcal P)\subseteq\mathcal L$.
Odwrotna inkluzja wynika z minimalności $\mathcal L$, ponieważ
$\sigma(\mathcal P)$ jest układem Dynkina zawierającym $\mathcal P$.
Ostatecznie $\mathcal L=\sigma(\mathcal P)\subseteq\mathcal D$.
\end{proof}

\begin{uwaga}
Twierdzenie $\pi$--$\lambda$ pozwala dowodzić własności wszystkich
zbiorów w $\sigma(\mathcal P)$, badając je najpierw na układzie
$\pi$. W rozdziale \ref{ch:miary-wlasnosci} zastosujemy je do
jednoznaczności miar; później również do jednoznaczności miar produktowych.
\end{uwaga}

\section{Klasy monotoniczne}

\begin{definicja}[Klasa monotoniczna]
\label{def:monotone-class}
Rodzinę $\mathcal M\subseteq\mathcal P(X)$ nazywamy \emph{klasą
monotoniczną}, jeżeli jest zamknięta na granice monotoniczne zbiorów:
\begin{align*}
A_n\in\mathcal M,\ A_n\uparrow A
&\Longrightarrow A\in\mathcal M,\\
B_n\in\mathcal M,\ B_n\downarrow B
&\Longrightarrow B\in\mathcal M.
\end{align*}
Zapis $A_n\uparrow A$ oznacza $A_n\subseteq A_{n+1}$ oraz
$A=\bigcup_nA_n$, a $B_n\downarrow B$ oznacza
$B_{n+1}\subseteq B_n$ oraz $B=\bigcap_nB_n$.
\end{definicja}

Przekrój klas monotonicznych jest klasą monotoniczną. Dlatego każda
rodzina $\mathcal E\subseteq\mathcal P(X)$ wyznacza najmniejszą klasę
monotoniczną $\mathcal M(\mathcal E)$, która ją zawiera.

\begin{twierdzenie}[Twierdzenie o klasie monotonicznej]
\label{thm:monotone-class}
Niech $X$ będzie zbiorem, a $\mathcal A$ algebrą podzbiorów $X$.
Wówczas $\mathcal M(\mathcal A)=\sigma(\mathcal A)$.
\end{twierdzenie}
\begin{proof}
Każda $\sigma$-algebra jest klasą monotoniczną, stąd
$\mathcal M(\mathcal A)\subseteq\sigma(\mathcal A)$.
Niech $\mathcal M=\mathcal M(\mathcal A)$. Pokażemy, że $\mathcal M$
jest $\sigma$-algebrą.

Najpierw rozważmy rodzinę
$\mathcal C=\{E\in\mathcal M:E^c\in\mathcal M\}$.
Jest ona klasą monotoniczną: dopełnienie zamienia ciągi rosnące
na malejące i odwrotnie. Ponieważ $\mathcal A$ jest algebrą,
$\mathcal A\subseteq\mathcal C$. Z minimalności $\mathcal M$
wynika $\mathcal C=\mathcal M$, czyli $\mathcal M$ jest zamknięta
na dopełnienia.

Ustalmy $A\in\mathcal A$ i połóżmy
$\mathcal M_A=\{E\in\mathcal M:E\cup A\in\mathcal M\}$.
Suma z ustalonym $A$ zachowuje granice monotoniczne, więc
$\mathcal M_A$ jest klasą monotoniczną. Algebraiczność $\mathcal A$
daje $\mathcal A\subseteq\mathcal M_A$, a zatem
$\mathcal M_A=\mathcal M$. Otrzymaliśmy
$E\cup A\in\mathcal M$ dla każdego $E\in\mathcal M$ i $A\in\mathcal A$.
Dla ustalonego $E\in\mathcal M$ rodzina
$\mathcal N_E=\{B\in\mathcal M:E\cup B\in\mathcal M\}$
jest więc klasą monotoniczną zawierającą $\mathcal A$, zatem
$\mathcal N_E=\mathcal M$. W konsekwencji $\mathcal M$ jest zamknięta
na sumy skończone.

Jeśli $E_1,E_2,\ldots\in\mathcal M$, to ich częściowe sumy
$F_n=\bigcup_{k=1}^nE_k$ należą do $\mathcal M$ i tworzą ciąg
rosnący. Stąd $\bigcup_nE_n\in\mathcal M$. Ponieważ
$X\in\mathcal A\subseteq\mathcal M$, otrzymaliśmy $\sigma$-algebrę.
Minimalność $\sigma(\mathcal A)$ daje drugą inkluzję.
\end{proof}

\begin{uwaga}[Dwa narzędzia rozszerzania własności]
Twierdzenia \ref{thm:pi-lambda} i \ref{thm:monotone-class} są podobne,
ale ich założenia nie są zamienne: twierdzenie Dynkina rozpoczyna się
od układu $\pi$, a twierdzenie o klasie monotonicznej --- od algebry.
Później użyjemy również wersji funkcyjnej twierdzenia o klasie
monotonicznej, po wprowadzeniu całki.
\end{uwaga}

\section{Zbiory borelowskie}

\begin{definicja}[$\sigma$-algebra borelowska]
\label{def:borel}
Niech $(X,\tau)$ będzie przestrzenią topologiczną, gdzie $\tau$
stanowi rodzinę jej zbiorów otwartych. Definiujemy
$\mathcal B(X)=\sigma_X(\tau)$.
Jest to \emph{$\sigma$-algebra borelowska}; jej elementy nazywamy
\emph{zbiorami borelowskimi}. Gdy $X=\mathbb R^d$, używamy standardowej
topologii euklidesowej.
\end{definicja}

Każdy zbiór otwarty i domknięty jest borelowski, podobnie jak każdy
przeliczalny przekrój zbiorów otwartych i przeliczalna suma zbiorów
domkniętych. Zbiory tych ostatnich dwóch typów oznacza się odpowiednio
$G_\delta$ i $F_\sigma$. Nie każdy zbiór borelowski ma jedną z tych
dwóch postaci.

\begin{twierdzenie}[Przeliczalne generatory na $\mathbb R$]
\label{thm:borel-generators-R}
Dla standardowej topologii na $\mathbb R$ zachodzą równości
\[
\mathcal B(\mathbb R)
=\sigma\bigl(\{(a,b):a,b\in\mathbb Q,\ a<b\}\bigr)
=\sigma\bigl(\{(-\infty,q):q\in\mathbb Q\}\bigr)
=\sigma\bigl(\{(-\infty,q]:q\in\mathbb Q\}\bigr).
\]
\end{twierdzenie}
\begin{proof}
Wszystkie wymienione przedziały są borelowskie, więc każda z
wygenerowanych rodzin jest zawarta w $\mathcal B(\mathbb R)$.
Każdy otwarty $U\subseteq\mathbb R$ jest sumą tych przedziałów
$(a,b)$ o końcach wymiernych, które zawierają się w $U$. Takich
przedziałów jest przeliczalnie wiele; daje to pierwszą równość.

Dalej, dla $a<b$ mamy $(a,b)=(-\infty,b)\setminus(-\infty,a]$ oraz
$(-\infty,a]=\bigcap_{n\geq1}(-\infty,a+1/n)$.
Wystarczy brać wymierne $a,b$ w pierwszym wzorze; w drugim
$a+1/n$ jest wymierne dla wymiernego $a$. To dowodzi drugiej
równości. Wreszcie dla $q\in\mathbb Q$ zachodzą
$(-\infty,q]=\bigcap_{n\geq1}(-\infty,q+1/n)$ oraz
$(-\infty,q)=\bigcup_{n\geq1}(-\infty,q-1/n]$,
co daje trzecią równość.
\end{proof}

\begin{twierdzenie}[Borelowskie prostokąty]
\label{thm:borel-Rd}
Dla $d\in\mathbb N$ $\sigma$-algebra $\mathcal B(\mathbb R^d)$
jest generowana przez otwarte prostokąty
$\prod_{j=1}^d(a_j,b_j)$ o wymiernych końcach. Ponadto
\[
\mathcal B(\mathbb R^d)=\mathcal B(\mathbb R)^{\otimes d},
\]
gdzie prawa strona oznacza $\sigma$-algebrę generowaną przez
prostokąty $B_1\times\cdots\times B_d$ z $B_j\in\mathcal B(\mathbb R)$.
\end{twierdzenie}
\begin{proof}
Otwarte prostokąty o wymiernych końcach tworzą przeliczalną bazę
topologii $\mathbb R^d$. Zatem dowolny zbiór otwarty jest sumą
przeliczalnej podrodziny tej bazy, co daje pierwsze twierdzenie.
Oznaczmy przez $\mathcal F$ prawą stronę deklarowanej równości.
Każdy prostokąt bazowy należy do $\mathcal F$, skąd
$\mathcal B(\mathbb R^d)\subseteq\mathcal F$.
Z drugiej strony rzuty współrzędnych $p_j:\mathbb R^d\to\mathbb R$
są ciągłe, więc $p_j^{-1}(B_j)$ jest borelowski dla
$B_j\in\mathcal B(\mathbb R)$: rodzina zbiorów $B$ o tej
własności jest $\sigma$-algebrą zawierającą otwarte podzbiory
$\mathbb R$. Ponieważ
$\prod_{j=1}^dB_j=\bigcap_{j=1}^dp_j^{-1}(B_j)$,
każdy prostokąt borelowski należy do $\mathcal B(\mathbb R^d)$.
Stąd $\mathcal F\subseteq\mathcal B(\mathbb R^d)$.
\end{proof}

Dla dowolnego zbioru $A\subseteq X$ przez $\chi_A$ oznaczamy
funkcję charakterystyczną: $\chi_A(x)=1$ dla $x\in A$
oraz $\chi_A(x)=0$ dla $x\notin A$.

\begin{przyklad}[Racjonalne i irracjonalne]
$\mathbb Q$ jest podzbiorem borelowskim $\mathbb R$, bo jest
przeliczalną sumą singletonów, a singletony są domknięte.
Zatem $\mathbb R\setminus\mathbb Q$ również jest borelowski.
Jednak funkcja charakterystyczna $\chi_{\mathbb Q}$ jest nieciągła
w każdym punkcie $\mathbb R$. Mierzalność nie jest więc bliska
ciągłości w sensie punktowym.
\end{przyklad}

\section{Odwzorowania mierzalne i ich charakteryzacje}

\begin{definicja}[Odwzorowanie mierzalne]
\label{def:measurable-map}
Niech $(X,\mathcal F)$ i $(Y,\mathcal G)$ będą przestrzeniami
mierzalnymi. Odwzorowanie $f:X\to Y$ nazywamy
\emph{$(\mathcal F,\mathcal G)$-mierzalnym}, jeżeli
$f^{-1}(B)\in\mathcal F$ dla każdego $B\in\mathcal G$.
Jeśli $Y$ jest przestrzenią topologiczną z
$\mathcal G=\mathcal B(Y)$, mówimy krótko, że $f$ jest
\emph{borelowsko mierzalne} względem $\mathcal F$.
\end{definicja}

\begin{twierdzenie}[Kryterium przez generatory]
\label{thm:measurable-generator}
Niech $(X,\mathcal F)$ będzie przestrzenią mierzalną, $Y$ zbiorem,
$\mathcal E\subseteq\mathcal P(Y)$ rodziną zbiorów i
$\mathcal G=\sigma_Y(\mathcal E)$. Wówczas odwzorowanie
$f:X\to Y$ jest $(\mathcal F,\mathcal G)$-mierzalne wtedy i tylko
wtedy, gdy $f^{-1}(E)\in\mathcal F$ dla wszystkich $E\in\mathcal E$.
\end{twierdzenie}
\begin{proof}
Konieczność wynika z inkluzji $\mathcal E\subseteq\mathcal G$.
Dla dostateczności niech
$\mathcal H=\{B\subseteq Y:f^{-1}(B)\in\mathcal F\}$.
To $\sigma$-algebra na $Y$, ponieważ obrazy odwrotne zachowują
dopełnienia i przeliczalne sumy. Z założenia
$\mathcal E\subseteq\mathcal H$, zatem
$\mathcal G=\sigma_Y(\mathcal E)\subseteq\mathcal H$.
\end{proof}

\begin{twierdzenie}[Złożenia i funkcje ciągłe]
\label{thm:composition-measurable}
Niech $(X,\mathcal F)$, $(Y,\mathcal G)$ i $(Z,\mathcal H)$
będą przestrzeniami mierzalnymi. Jeżeli $f:X\to Y$ jest
$(\mathcal F,\mathcal G)$-mierzalne, a $g:Y\to Z$ jest
$(\mathcal G,\mathcal H)$-mierzalne, to $g\circ f$ jest
$(\mathcal F,\mathcal H)$-mierzalne.
Ponadto każde odwzorowanie ciągłe $g:Y\to Z$ pomiędzy
przestrzeniami topologicznymi jest
$(\mathcal B(Y),\mathcal B(Z))$-mierzalne.
\end{twierdzenie}
\begin{proof}
Dla $C\in\mathcal H$ mamy
$(g\circ f)^{-1}(C)=f^{-1}(g^{-1}(C))\in\mathcal F$.
Jeśli $g$ jest ciągłe, obrazy odwrotne zbiorów otwartych są otwarte;
z kryterium \ref{thm:measurable-generator} wynika jego mierzalność
borelowska.
\end{proof}

\begin{twierdzenie}[Funkcje rzeczywiste i działania]
\label{thm:measurable-real}
Niech $(X,\mathcal F)$ będzie przestrzenią mierzalną.
Dla funkcji $f:X\to\mathbb R$ następujące warunki są równoważne:
\begin{enumerate}
\item $f$ jest $(\mathcal F,\mathcal B(\mathbb R))$-mierzalna;
\item $\{x\in X:f(x)<q\}\in\mathcal F$ dla każdego $q\in\mathbb Q$;
\item $\{x\in X:f(x)\leq q\}\in\mathcal F$ dla każdego $q\in\mathbb Q$.
\end{enumerate}
Jeżeli $f,g:X\to\mathbb R$ są mierzalne, to funkcje
$f+g$, $fg$, $\max(f,g)$ i $\min(f,g)$ są mierzalne.
\end{twierdzenie}
\begin{proof}
Równoważność trzech warunków wynika z twierdzeń
\ref{thm:borel-generators-R} i \ref{thm:measurable-generator}.
Odwzorowanie $x\mapsto(f(x),g(x))$ do $\mathbb R^2$ jest mierzalne:
obrazy odwrotne prostokątów $B_1\times B_2$, gdzie
$B_1,B_2\in\mathcal B(\mathbb R)$, są przekrojami
$f^{-1}(B_1)\cap g^{-1}(B_2)$. Prostokąty te generują
$\mathcal B(\mathbb R^2)$ na mocy twierdzenia
\ref{thm:borel-Rd}. Funkcje na $\mathbb R^2$ określone przez
$(u,v)\mapsto u+v$, $(u,v)\mapsto uv$,
$(u,v)\mapsto\max(u,v)$ i $(u,v)\mapsto\min(u,v)$ są ciągłe.
Teza wynika więc z twierdzenia \ref{thm:composition-measurable}.
\end{proof}

\begin{przyklad}[Funkcje charakterystyczne zbiorów]
Dla $A\subseteq X$ definiujemy funkcję charakterystyczną
$\chi_A:X\to\mathbb R$ wzorem
$\chi_A(x)=1$ dla $x\in A$ i $\chi_A(x)=0$ dla $x\notin A$.
Wtedy $\chi_A$ jest mierzalna względem $\mathcal F$
dokładnie wtedy, gdy $A\in\mathcal F$.
Dostateczność wynika z tego, że funkcja przyjmuje tylko dwie
wartości; konieczność z równości
$A=\chi_A^{-1}((1/2,\infty))$.
\end{przyklad}

\begin{przyklad}[Mierzalność nie implikuje regularności geometrycznej]
Na $X=\mathbb R$ z $\mathcal F=\mathcal B(\mathbb R)$ funkcja
$\chi_{\mathbb Q}$ jest mierzalna, choć nie jest ciągła
w żadnym punkcie. Z drugiej strony, jeśli na nieprzeliczalnym
zbiorze $X$ przyjmiemy $\mathcal F=\{\varnothing,X\}$, to
każda mierzalna funkcja $X\to\mathbb R$ jest stała. Rzeczywiście,
gdyby przyjmowała wartości $a<b$, to dla pewnej liczby wymiernej
$q\in(a,b)$ zbiór $\{f<q\}$ byłby niepusty i różny od $X$.
\end{przyklad}

\begin{uwaga}[Co zostało, a co nie zostało wykazane]
Wprowadziliśmy pojęcie mierzalności bez definiowania jakiejkolwiek
miary. To rozróżnienie jest zasadnicze: przestrzeń mierzalna określa,
\emph{którym} zbiorom wolno przypisać rozmiar; miara określa
\emph{jaki} jest ten rozmiar. Miary pojawią się w rozdziale
\ref{ch:miary-wlasnosci}, a mierzalność funkcji o wartościach
w $[-\infty,\infty]$ omówimy szczegółowo przed zdefiniowaniem całki.
\end{uwaga}